%=====================================================================
\chapter{Glossary}\label{app:glossary}
%=====================================================================
\normalfigures

Terms are defined as this course uses them. Where a word is used differently
elsewhere in the industry, that is said. The chapter in brackets is where the
term is introduced and where the evidence for any claim about it lives.

\begin{multicols}{2}
\small
\raggedright

\textbf{abandonment.} A basket or checkout the customer began and did not
finish. In this platform it is visible as the gap between rows in
\texttt{oc\_order} and orders above status 0. (\chref{checkout},
\chref{dashboard})

\textbf{acquirer.} The bank holding the shop's account, which accepts card
payments on its behalf and carries the risk that the shop fails to deliver.
(\chref{payments})

\textbf{affiliate.} A customer paid commission on orders arriving through
their tracking code. Default commission here is 5\%. (\chref{discounts})

\textbf{authorisation.} A request asking the issuer whether a payment would
succeed, reserving the amount. Reversible, instant, and not money.
(\chref{payments})

\textbf{bcrypt.} A deliberately slow password hashing function with a tunable
cost factor. PHP's \texttt{PASSWORD\_DEFAULT} produces bcrypt at cost 10 here.
(\chref{accounts})

\textbf{B2B, B2C, C2C.} Business to business, to consumer, and consumer to
consumer. A statement about who the parties are, not about the software.
(\chref{models})

\textbf{capture.} Claiming an amount previously authorised. Only now is the
shop owed anything. (\chref{payments})

\textbf{cash on delivery (COD).} Payment collected by the courier on handover.
Dominant in this market, and the most expensive instrument once refusals are
counted. (\chref{payments})

\textbf{chargeback.} A card payment reversed by the issuer after a dispute.
The loss falls on the shop. (\chref{payments})

\textbf{container.} A running instance of an image, with its own filesystem and
process space, sharing the host's kernel. Not a virtual machine.
(\chref{stack})

\textbf{content security policy (CSP).} A response header listing where a page
may load scripts, styles and images from. Enforced by the browser.
(\chref{security})

\textbf{cookie.} A named value the server asks the browser to store and return
on later requests. The only standard way to relate two requests. Lives on the
customer's machine, where they can read and change it. (\chref{basket})

\textbf{cost factor.} Bcrypt's work parameter. Each step doubles the time for
you and for an attacker. (\chref{accounts})

\textbf{coupon.} A discount the customer must type a code to obtain. The code
is what makes a campaign measurable. (\chref{discounts})

\textbf{credential stuffing.} Trying email and password pairs stolen elsewhere.
Defeats password rules, because the passwords are correct. (\chref{accounts})

\textbf{CSRF token.} A one-time value in a form, checked on submission, so that
another site cannot submit the form on a visitor's behalf. This shop's login
carries one. (\chref{accounts}, \chref{security})

\textbf{destination principle.} Tax belongs where the goods are consumed ---
hence a tax rule based on the shipping address. (\chref{shipping})

\textbf{Docker Compose.} A file describing several containers, their networks
and their volumes, brought up together by one command. (\chref{stack})

\textbf{gateway (payment).} The software company between the shop and the
acquirer. The only one of the five parties you integrate with.
(\chref{payments}, \chref{gateway})

\textbf{geo zone.} A named set of (country, zone) pairs that shipping methods
and tax rules attach to. It has no meaning of its own. (\chref{shipping})

\textbf{hosted page.} A payment architecture where the customer types their
card on the provider's site. Keeps the shop in PCI DSS SAQ-A.
(\chref{gateway})

\textbf{idempotency key.} A value sent with a request so that repeating it has
the effect of sending it once. Must belong to the order, not the request.
(\chref{gateway})

\textbf{image (container).} A read-only layered filesystem holding an
application and its dependencies. Built once, identical everywhere.
(\chref{stack})

\textbf{issuer.} The bank that gave the customer their card or holds their
wallet. Decides whether to approve, and hears their complaints.
(\chref{payments})

\textbf{long tail.} A latency distribution where a minority of requests take
far longer than the median, so the mean sits well above it. A timeout must be
chosen against the tail. (\chref{gateway})

\textbf{marketplace.} A venue where third parties sell, with the operator
taking a fee and holding neither stock nor the customer relationship.
(\chref{models})

\textbf{merchant discount rate.} The single percentage a shop pays on a card
payment, divided afterwards between scheme, acquirer and gateway.
(\chref{payments})

\textbf{module.} A unit of functionality added to the platform without editing
it: a shipping method, a payment method, a dashboard panel.
(\chref{modules})

\textbf{option (product).} A choice offered on a product --- size, colour ---
stored as rows that grow by addition. Contrast \emph{variant}.
(\chref{variations})

\textbf{order.} A row recording a commitment to buy, with the products,
totals and addresses \emph{copied} as they were at the time. It does not point
at the catalogue. (\chref{checkout})

\textbf{order history.} The list of an order's status changes with dates. The
status says where an order is; the history says how it got there, and is what
a dispute is settled from. (\chref{checkout}, \chref{orders})

\textbf{order status.} The state an order is in. The set is the shop's own, not
a standard. Status 0 is not in the set. (\chref{checkout})

\textbf{PayPak.} Pakistan's domestic card scheme, running over the 1LINK
switch. Cheaper to accept and domestic only. (\chref{payments})

\textbf{PCI DSS.} The card industry's data security standard. Its reporting
burden is graded by how much card data you touch; a hosted page keeps you on
the shortest questionnaire. (\chref{gateway})

\textbf{PECA 2016.} Pakistan's Prevention of Electronic Crimes Act.
Unauthorised access to an information system is an offence under it, with no
educational exemption. (\chref{security}, \chref{legal})

\textbf{percentile.} The value below which a given share of measurements fall.
The 95th percentile of a response time tells you what a slow customer
experiences; the mean does not. (\chref{gateway})

\textbf{Raast.} The State Bank's instant payment system, settling bank to bank
in seconds at negligible cost. (\chref{payments})

\textbf{return-to-origin.} A delivery refused at the door and sent back. The
shop pays both legs and keeps the goods. (\chref{payments})

\textbf{salt.} A random per-user value mixed into a password hash, so that
identical passwords hash differently and precomputed tables are useless. Need
not be secret. (\chref{accounts})

\textbf{session.} Server-side state filed under a key, with the browser
carrying only the key. Distinct from the cookie that points at it.
(\chref{basket})

\textbf{session engine.} Where session records are kept --- database, files or
Redis. It does not decide where the basket is. (\chref{basket})

\textbf{settlement.} The acquirer actually paying the shop, days later and net
of fees. (\chref{payments})

\textbf{sort order (totals).} The position of a total line in the pipeline. It
decides what each line is computed from, and therefore what the customer pays.
(\chref{discounts})

\textbf{special.} A dated replacement price for a product and customer group,
applied automatically. (\chref{discounts})

\textbf{tax base.} Which address a tax rule is decided by: store, payment or
shipping. The most consequential dropdown in the admin panel.
(\chref{shipping})

\textbf{tax class.} A kind of product for tax purposes, such as taxable goods
or downloadable products. (\chref{shipping})

\textbf{tax rate.} A number and the geo zone it applies in. Percentage or fixed
amount. (\chref{shipping})

\textbf{tax rule.} The join between a class and a rate, with a priority and a
base. (\chref{shipping})

\textbf{template override.} A copy of one theme file in a child theme,
replacing the original for that one page. The override directory is the list of
your changes. (\chref{ux})

\textbf{theme.} A directory of Twig templates and assets. Never edit the
default one. (\chref{ux})

\textbf{timing side channel.} Information leaked by how long an operation
takes. Here, a registered email answered 65\,ms slower than an unknown one.
(\chref{accounts})

\textbf{Twig.} The template language this platform renders pages with. It
auto-escapes output, which is why the script probes in \chref{security} did
nothing.

\textbf{user enumeration.} Discovering which accounts exist, by timing,
messages or behaviour. Reconnaissance rather than an attack, and worth closing
for that reason. (\chref{accounts})

\textbf{variant.} A product generated per combination of choices, growing by
multiplication. Contrast \emph{option}. (\chref{variations})

\textbf{volume (container).} Storage managed by the container runtime and
outliving any container. \texttt{down} keeps it; \texttt{down -v} deletes it.
(\chref{stack}, \chref{basket})

\textbf{voucher (gift).} A stored-value instrument bought by one person and
spent by another. HEC's outline names it; \textbf{this platform does not have
one}. (\chref{discounts})

\end{multicols}
